% !TEX root = ../frode-bibelen.tex
% ----------------------------------------------------------------------------
%  Frode-bibelen -- 1 Eksamen: format og strategi
%  Inkluderes fra frode-bibelen.tex (kompiler alltid hovedfilen).
% ----------------------------------------------------------------------------

% ======================================================== 1 EKSAMEN ========
\kapittel{Eksamen: format og strategi}

Alle eksamenene fra H19 til S26 har hatt samme form: \textbf{4 timer}, digitalt i Inspera,
\textbf{alle trykte og skrevne hjelpemidler} (ikke kalkulator), og to uavhengige oppgaver.

\begin{center}
\begin{tabular}{@{}L{2.3cm}L{1.2cm}L{11.8cm}@{}}
\toprule
\textbf{Oppgave} & \textbf{Vekt} & \textbf{Innhold} \\
\midrule
1a & 30\,\% samlet & \emph{«Hva blir utskriften?»} -- tekst/\cc{char}-array, pekere og \cc{string.h}-funksjoner (\cc{strchr}, \cc{strstr}, \cc{strncmp} \dots). Antall utskriftslinjer oppgis som hjelp. \\
1b & & \emph{«Hva blir utskriften?»} -- \cc{struct}, funksjoner med parametre (kopi vs.\ peker), \cc{bool}, \cc{strcpy}/\cc{strcat}. \\
2 (A--G/H) & 70\,\% & Et større menystyrt program. Et \textbf{skjelett} (\cc{.tpl}) ligger som vedlegg i Inspera: \cc{main}, \cc{skrivMeny}, \cc{struct}-er, prototyper og globale variable er som regel ferdige. Du skriver innmaten i funksjonene. \\
\bottomrule
\end{tabular}
\end{center}

\subsection{Forslag til tidsbruk (240 minutter)}
\begin{tabular}{@{}L{3.2cm}L{12.5cm}@{}}
0--10 min & Les \emph{hele} oppgavesettet. Åpne skjelettet. Marker hva som er ferdiglaget. \\
10--75 min & Oppgave 1a og 1b (30\,\% $\approx$ 70 min). Bruk sporingsarket (kapittel~\ref{kap:sporingsark}). \\
75--225 min & Oppgave 2. Start med deloppgavene du er tryggest på (ofte A, B og G). \\
225--240 min & Gjennomgang med sjekklistene (kapittel~\ref{kap:sjekk1} og~\ref{kap:sjekk2}). \\
\end{tabular}

\subsection{Frodes eksamensråd (fritt gjengitt)}
Kilde: \href{https://frh.folk.ntnu.no/grprog/ex_raad.html}{\cc{grprog/ex\_raad.html}}.
\begin{frode}[{Hva oppgave 2 egentlig tester}]
\begin{itemize}
  \item Det viktigste er å vise at du kan \textbf{opprette, finne fram i og behandle datastrukturer}.
  \item \textbf{Les hele oppgaveteksten nøye} før du skriver. Tegn datastrukturen og lag notater/pseudokode -- både for helheten og for hver funksjon.
  \item \textbf{Ikke la noe stå blankt.} Prøv å få «minst en halv F» på alle deloppgaver. Skriv heller ryddig pseudokode enn ingenting.
  \item \textbf{Ikke gjenta oppgaveteksten} eller læreboka -- henvis heller.
  \item \textbf{Kommenter det ikke-opplagte.} Er du usikker på syntaksen, skriv i en kommentar hva du \emph{vil} gjøre. Riktig tankegang teller selv om koden har feil.
  \item \textbf{Robusthet vektlegges ikke} med mindre oppgaven ber om det: anta at input er korrekt og at filene følger formatet.
  \item Sensor ser \textbf{ikke} etter manglende \cc{;} eller parenteser, men om \textbf{hver funksjon gjør jobben sin}.
  \item Løsningsforslagene er «overløst». Du kan få god karakter uten like elegant eller helt syntaktisk korrekt kode.
  \item «Gullkornet»: Ingenting \emph{nytt} kommer på eksamen -- det du kan, skal bare kombineres på en annen måte.
  \item Ta en pause underveis, og slapp av dagen før.
\end{itemize}
\end{frode}

\subsection{Oppgave 2 sin form}
\begin{itemize}
  \item \textbf{Datastrukturen} er helt eller delvis gitt og forklart. Nesten alltid: en global \textbf{array med pekere til \cc{struct}} pluss en teller (\cc{gKunder[]}, \cc{gAntallKunder}). Noen ganger to slike (f.eks.\ kunder og banker), eller en fast array av \cc{struct}-er.
  \item \textbf{Variable, \cc{enum}, \cc{struct} og funksjoner}: alt eller det meste er gitt. Noen ganger må du lage en hjelpefunksjon selv (merket «NY» eller «EKSTRA» i løsningsforslagene).
  \item \textbf{Løsningsforslag og rettemal} legges i \href{https://frh.folk.ntnu.no/grprog/eksamen/}{eksamensarkivet} ca.\ en time etter eksamen.
\end{itemize}

\subsection{Typiske deloppgaver i oppgave 2 -- og hvor malen står}
\begin{center}
\begin{tabular}{@{}L{5.6cm}L{7.6cm}L{2.0cm}@{}}
\toprule
\textbf{Deloppgaven ber deg \dots} & \textbf{Funksjonspar (eksempel)} & \textbf{Mal} \\
\midrule
skrive alle (hoveddata) & \cc{skrivAlleKunder} + \cc{kundeSkrivData} & \ref{sec:skrivalle} \\
legge inn en ny & \cc{nyKunde} + \cc{kundeLesData} & \ref{sec:ny}, \ref{sec:lesdata} \\
legge inn ny med unikt navn & \cc{finnKunde} (indeks / $-1$) & \ref{sec:unik} \\
skrive \emph{alt} om én valgt & \cc{skrivEnKunde} + \cc{kundeSkrivAlt} & \ref{sec:velgen} \\
endre / registrere noe på én & \cc{endreKunde} + \cc{kundeEndre} & \ref{sec:endre} \\
fjerne / slette én & \cc{fjernKunde} + \cc{kundeFrigi} & \ref{sec:slett} \\
skrive alle som matcher (type, tekst) & \cc{kundeErType}, \cc{strstr} & \ref{sec:match} \\
finne flest / størst / dag med mest & hjelpearray + maks-søk & \ref{sec:statistikk} \\
\cc{enum} inn og ut & \cc{lesType}, \cc{skrivType} & \ref{sec:enum} \\
skrive alt til fil & \cc{skrivTilFil} + \cc{kundeSkrivTilFil} & \ref{sec:skrivfil} \\
lese alt fra fil (oftest G) & \cc{lesFraFil} + \cc{kundeLesFraFil} & \ref{sec:lesfil} \\
\bottomrule
\end{tabular}
\end{center}

\subsection*{Før eksamensdagen}
\boks Bibelen skrevet ut og satt i perm \quad \boks \cc{LesData.h} skrevet ut \quad \boks 2--3 løsningsforslag skrevet ut\\
\boks Prøvd digital eksamen i Inspera \quad \boks Løst minst tre gamle eksamener på tid \quad \boks Sovet
